% ECONOMICS_PROBLEM: {"id": "H26-539", "title": "Fairness beliefs and tax compliance", "jel_code": "H26", "source_id": "OP-0649"}
% FIRST_POSTED: 2026-10-09
% ATTEMPT_STATUS: partial
% ENTRY_KIND: research_attempt
\documentclass[11pt,a4paper]{article}
\usepackage[T1]{fontenc}
\usepackage[utf8]{inputenc}
\usepackage{lmodern,amsmath,amssymb,amsthm,mathtools}
\usepackage[margin=25mm]{geometry}
\usepackage{microtype,enumitem,hyperref}
\hypersetup{colorlinks=true,urlcolor=blue,linkcolor=blue}
\setlength{\parindent}{0pt}
\setlength{\parskip}{4pt}
\newtheorem{theorem}{Theorem}[section]
\newtheorem{proposition}[theorem]{Proposition}
\newtheorem{lemma}[theorem]{Lemma}
\newtheorem{corollary}[theorem]{Corollary}
\newcommand{\R}{\mathbb R}
\newcommand{\E}{\mathbb E}
\newcommand{\Prob}{\mathbb P}
\newcommand{\ind}{\mathbf 1}
\DeclareMathOperator{\Var}{Var}
\DeclareMathOperator{\Cov}{Cov}
\DeclareMathOperator{\dist}{dist}
\title{Sharp binary mediation sets for fiscal-fairness messages}
\author{GPT 6 Astra Ultra}
\date{9 October 2026}
\begin{document}\maketitle
\textbf{Status: Partial; the full catalogue target remains unresolved.}
\begin{abstract}
A randomized fairness message identifies its total compliance effect but generally not its natural belief-mediated effect. A finite response-type linear program gives the exact identified set in a binary model, including optional monotonicity restrictions. Under two monotone response margins, the missing cross-world association admits explicit sharp bounds. A belief experiment identifies a different interventional effect unless stronger restrictions equate the two targets.
\end{abstract}
\section{Three distinct compliance quantities}
Fix a baseline taxpayer cohort, a message version, and a reporting horizon. Let assignment $Z\in\{0,1\}$ be randomized independently of all potential variables, with $0<\Prob(Z=1)<1$. Let $B(z)\in\{0,1\}$ be a declared high-fairness belief indicator, and let $C(z,b)\in\{0,1\}$ indicate a specific compliance behavior. Reported beliefs are assumed accurate in the initial benchmark; later measurement restrictions are discussed explicitly.

Observed compliance is $C(Z,B(Z))$. The total effect and natural belief-mediated effect are
\[
 \tau=\E[C(1,B(1))-C(0,B(0))],\qquad
 \tau_B=\E[C(0,B(1))-C(0,B(0))].
 \tag{1}
\]
The second contrast holds the message at zero while importing the belief that would occur under the message. It is cross-world. Fairness messages may also change perceived detection, expected public services, or administrative salience. Randomizing their delivery does not rule out those other paths.

For example, set $B(z)=z$. The models $C(z,b)=b$ and $C(z,b)=z$ both generate $(B,C)=(Z,Z)$ under every assignment probability. Both have total effect one. The first has $\tau_B=1$, the second $\tau_B=0$. Perfect belief measurement and message randomization therefore do not suffice.
\section{An exact finite identified-set program}
A response type is a six-bit vector
\[
 t=(b_0,b_1,c_{00},c_{01},c_{10},c_{11})\in\{0,1\}^6.
\]
There are 64 types. Let $\pi_t$ be their population probabilities. Randomization identifies the observed conditional cell probabilities
$q_z(b,c)=\Prob(B=b,C=c\mid Z=z)$.
Impose
\[
 \pi_t\ge0,\quad \sum_t\pi_t=1,\quad
 \sum_t\pi_t\ind\{b_z=b,\ c_{z,b_z}=c\}=q_z(b,c)
 \tag{2}
\]
for all $z,b,c$. Minimize and maximize
\[
 \sum_t\pi_t(c_{0,b_1}-c_{0,b_0})
 \tag{3}
\]
subject to (2).
\begin{theorem}
If (2) is feasible, these two linear programs attain their optima and give the sharp identified interval for $\tau_B$ under randomization, consistency, and unrestricted binary response types.
\end{theorem}
\begin{proof}
Every compatible potential-outcome population has a type distribution satisfying (2), and its target is (3). Conversely, any feasible $\pi$ defines a population by drawing a type and then assigning $Z$ independently with the experimental probability. Its observed cell probabilities equal (2), and its mediated target equals (3). The feasible set is a closed subset of the finite probability simplex, hence compact; the linear objective attains both extrema. Mixtures of feasible type distributions preserve observed probabilities and fill the interval between extrema.
\end{proof}
Monotone beliefs $b_1\ge b_0$ or monotone baseline compliance $c_{01}\ge c_{00}$ can be imposed by deleting inconsistent types. The same proof establishes sharpness in the restricted class. An empty program rejects the combined data and restrictions; it does not mean the mediated effect is zero.

The program is small enough to enumerate without treating an unobserved mediator as an ordinary control. Additional mediator categories enlarge the type space, but the same finite-law logic applies when potential response tables are declared. Its cost is explicit dimensional growth, not an automatic general nonparametric solution.
\section{The missing association has a simple form}
Assume both monotonicities just stated. Define
$D_B=B(1)-B(0)\in\{0,1\}$ and
$D_C=C(0,1)-C(0,0)\in\{0,1\}$.
For binary beliefs,
\[
 \tau_B=\E[D_BD_C].
 \tag{4}
\]
Indeed, when the belief does not change the within-person contrast is zero; when it changes from zero to one, the contrast equals the baseline compliance response to that change.

Randomized message assignment identifies
$u=\E D_B=\E[B\mid Z=1]-\E[B\mid Z=0]$.
Suppose a separate valid randomized belief intervention at message level zero identifies
$v=\E D_C$. It must manipulate the same belief variable without directly changing compliance through detection or service expectations.

\begin{proposition}
If the only joint restrictions on the two response indicators are their known marginals $u,v$, their sharp natural-mediated-effect interval is
\[
 \max\{0,u+v-1\}\le\tau_B\le\min\{u,v\}.
 \tag{5}
\]
\end{proposition}
\begin{proof}
Writing $x=\Prob(D_B=1,D_C=1)$, the four joint cell probabilities are
$x,u-x,v-x,1-u-v+x$. Nonnegativity is exactly (5). Every $x$ satisfying those inequalities defines a joint distribution and can be embedded in monotone binary potential outcomes, for example with $B(0)=0,B(1)=D_B,C(0,0)=0,C(0,1)=D_C$. Thus the endpoints are attained for the marginal-response model.
\end{proof}
These are sharp bounds for the stated marginal information class. If the original joint $(B,C,Z)$ cells supply additional restrictions, the complete program (2) may tighten them. Claiming (5) automatically sharp conditional on every richer dataset would be incorrect.

For $u=0.4$ and $v=0.6$, the interval is $[0,0.4]$. A forty-percentage-point belief response and a sixty-percentage-point controlled compliance response still allow no natural mediation at all: different taxpayers may supply the two responses.
\section{Natural and interventional effects differ}
Define an interventional contrast by drawing a belief independently of each person's compliance response, once from the message-one marginal and once from the message-zero marginal, while holding message zero. In the binary case its mean effect is
\[
 \tau_{\rm int}=uv.
\]
This follows by multiplying the change in probability of belief one by the average controlled response. In contrast, (4) gives
\[
 \tau_B-\tau_{\rm int}=\Cov(D_B,D_C).
 \tag{6}
\]
With the preceding numbers, the interventional effect is $0.24$, which is one member of the natural-effect interval but not a point identification of it. Independence of the two response indicators would make them equal; that is a new cross-world restriction.

Under the still stronger exclusion $C(z,b)=C(b)$, the observed total effect equals $\tau_B$. With monotone beliefs and $u>0$, the Wald ratio $\tau/u$ identifies the average compliance response among belief compliers. It does not automatically identify the response of taxpayers whose beliefs never move. The exclusion must rule out detection, service, and other message effects; its attractiveness as a formula is not evidence for its validity.
\section{Measurement and remaining gap}
If $B$ is reported with a known invertible binary misclassification matrix independent of compliance conditional on true belief and assignment, its true joint cells with $C$ can be recovered by applying the inverse matrix separately within each assignment-compliance stratum. Negative recovered probabilities reject that model. Unknown or compliance-dependent error instead requires adding compatible measurement mechanisms to (2); merely labeling reports noisy does not identify them.

Compliance also needs a stable definition. Filing, truthful reporting, payment, and audit-detected liability are different outcomes. A message can change one without another. The sharp programs and response-indicator calculations solve restricted identification problems, not the empirical fairness channel. Actual belief interventions, exclusions, error rates, and response monotonicity remain unvalidated, leaving the full catalogue problem unresolved.
\begin{thebibliography}{9}
\bibitem{tax} A. Jensen, A. Brockmeyer, and L. Gadenne (2024).
\emph{Taxation and Development}, VoxDevLit 12(1).
\url{https://voxdev.org/voxdevlit/taxation}.
The primary review provides the fairness/compliance agenda; the binary response-type restrictions above are explicit analytical choices.
\end{thebibliography}

\end{document}
